\documentclass[11pt]{article}

\usepackage{amsmath,amssymb,amsthm}

\newtheorem{theorem}{Theorem}
\newtheorem{lemma}[theorem]{Lemma}
\theoremstyle{remark}
\newtheorem*{remark}{Remark}

\newcommand{\F}{\mathbb F}
\newcommand{\zetaP}{\zeta_p}
\newcommand{\card}[1]{\lvert #1\rvert}
\newcommand{\ip}[2]{\langle #1,#2\rangle}

\title{The Depth of a Bent Partition Is a Power of the Characteristic}
\author{}
\date{}

\begin{document}

\maketitle

\begin{abstract}
Let $p$ be a prime and let $n>0$ be even.  We prove, without any
regularity or structural hypothesis, that the depth of every bent
partition of $\F_p^n$ is a power of $p$.  The proof averages, over all
balanced labelings of the blocks, the number of zero values of one fixed
nonzero derivative.
\end{abstract}

\section{Definitions and exact statement}

Put $V=\F_p^n$, equip $V$ with the standard $\F_p$-valued inner
product, and set
\[
  \zetaP=\exp(2\pi i/p).
\]
For a function $f:V\to\F_p$, its Walsh transform is
\[
  W_f(b)=\sum_{x\in V}\zetaP^{\,f(x)-\ip{b}{x}}
  \qquad (b\in V).
\]
The function $f$ is \emph{bent} if
\[
  \card{W_f(b)}=p^{n/2}
  \qquad\text{for every }b\in V.
\]

Let $\Omega=\{A_1,\ldots,A_K\}$ be a partition of $V$ into nonempty
sets, and suppose that $p\mid K$.  A labeling
\(\ell:\{1,\ldots,K\}\to\F_p\) is called \emph{balanced} if
\[
  \card{\ell^{-1}(c)}=K/p
  \qquad\text{for every }c\in\F_p.
\]
It induces the function $f_\ell:V\to\F_p$ defined by
\[
  f_\ell(x)=\ell(i)\qquad\text{when }x\in A_i.
\]
The partition $\Omega$ is a \emph{bent partition of depth $K$} if
$f_\ell$ is bent for every balanced labeling $\ell$.

\begin{theorem}[Exact depth theorem]\label{thm:main}
Let $p$ be a prime, let $n>0$ be even, and let
$\Omega=\{A_1,\ldots,A_K\}$ be a bent partition of $\F_p^n$ of depth
$K$ in the sense defined above.  Then there is an integer $k$ with
\[
  1\leq k\leq n
\]
such that
\[
  K=p^k.
\]
No hypothesis is imposed on regularity, weak regularity, the sizes or
types of the blocks, normality, spread structure, vectorial-dual-bent
structure, or the manner in which the partition was constructed.
\end{theorem}

\section{A derivative lemma}

For $a\in V$, define
\[
  D_a f(x)=f(x+a)-f(x)
\]
and the autocorrelation
\[
  C_f(a)=\sum_{x\in V}\zetaP^{\,D_a f(x)}.
\]

\begin{lemma}\label{lem:derivative}
If $f:V\to\F_p$ is bent and $a\ne0$, then $D_a f$ is balanced.  More
precisely, for every $c\in\F_p$,
\[
  \#\{x\in V:D_a f(x)=c\}=p^{n-1}.
\]
\end{lemma}

\begin{proof}
Writing $a=x-y$ in the product $W_f(b)\overline{W_f(b)}$ gives the
finite Fourier identity
\begin{equation}\label{eq:wiener-khinchin}
  \card{W_f(b)}^2
  =\sum_{a\in V} C_f(a)\zetaP^{-\ip{b}{a}}.
\end{equation}
Since $f$ is bent, the left-hand side of
\eqref{eq:wiener-khinchin} is the constant $p^n$.  Fourier inversion on
the additive group of $V$ therefore gives
\[
  C_f(a)=0\qquad(a\ne0).
\]

Fix $a\ne0$, and for $c\in\F_p$ put
\[
  m_c=\#\{x\in V:D_a f(x)=c\}.
\]
After identifying $\F_p$ with $\{0,1,\ldots,p-1\}$ in exponents, the
preceding equality says
\[
  \sum_{c=0}^{p-1}m_c\zetaP^c=0.
\]
Thus the polynomial $P(X)=\sum_{c=0}^{p-1}m_cX^c\in\mathbb Z[X]$
vanishes at $\zetaP$.  Because $p$ is prime, the minimal polynomial of
$\zetaP$ over $\mathbb Q$ is
\[
  \Phi_p(X)=1+X+\cdots+X^{p-1}.
\]
As $\deg P\le p-1=\deg\Phi_p$, it follows that
$P(X)=t\Phi_p(X)$ for some rational number $t$.  Hence all the $m_c$
are equal.  Their sum is $\card V=p^n$, so
$m_c=p^n/p=p^{n-1}$ for every $c\in\F_p$.
\end{proof}

\section{A balanced-labeling lemma}

\begin{lemma}\label{lem:balanced-labeling}
Let $K=pr$, and choose uniformly at random a labeling
$\ell:\{1,\ldots,K\}\to\F_p$ satisfying
\[
  \card{\ell^{-1}(c)}=r
  \qquad\text{for every }c\in\F_p.
\]
For any two distinct block indices $i\ne j$,
\[
  \Pr\bigl(\ell(i)=\ell(j)\bigr)=\frac{r-1}{K-1}.
\]
This probability is independent of the sizes of the blocks and of the
choice of the distinct indices $i,j$; it depends only on $K,p,r$.
\end{lemma}

\begin{proof}
For every $c\in\F_p$, uniformity of the balanced labelings gives
\[
  \Pr\bigl(\ell(i)=c\bigr)=\frac{r}{K}=\frac1p.
\]
Conditional on $\ell(i)=c$, exactly $r-1$ of the remaining $K-1$
block indices have label $c$.  Hence
\[
\begin{aligned}
  \Pr\bigl(\ell(i)=\ell(j)\bigr)
  &=\sum_{c\in\F_p}
    \Pr\bigl(\ell(i)=c\bigr)
    \Pr\bigl(\ell(j)=c\mid\ell(i)=c\bigr)\\
  &=p\cdot\frac1p\cdot\frac{r-1}{K-1}
   =\frac{r-1}{K-1}.
\end{aligned}
\]
The calculation uses only the indexed set of $K$ blocks and their label
multiplicities, not their cardinalities, and it is identical for every
ordered pair of distinct indices $i,j$.
\end{proof}

\section{Proof of the depth theorem}

\begin{proof}[Proof of Theorem~\ref{thm:main}]
Write
\[
  Q=p^n,\qquad r=K/p.
\]
Here $r$ is a positive integer.  Since $n>0$, choose and fix
$a\in V\setminus\{0\}$.  Let
\begin{equation}\label{eq:T-def}
  T_a=
  \#\bigl\{x\in V:\text{$x$ and $x+a$ belong to the same block of
  $\Omega$}\bigr\}.
\end{equation}
In particular, $T_a$ is a nonnegative integer.

Choose a balanced labeling $\ell$ uniformly at random from the finite
set of all balanced labelings.  For every such $\ell$, Lemma
\ref{lem:derivative}, applied to the bent function $f_\ell$, shows that
\begin{equation}\label{eq:zero-count}
  \#\{x\in V:D_a f_\ell(x)=0\}=Q/p.
\end{equation}

We now compute the expectation of the left-hand side of
\eqref{eq:zero-count} by considering each $x\in V$.  If $x$ and $x+a$
belong to the same block, then
$D_a f_\ell(x)=0$ for every labeling $\ell$.  If they belong to two
distinct blocks, then $D_a f_\ell(x)=0$ exactly when those two blocks
receive the same label.  By Lemma~\ref{lem:balanced-labeling},
\[
  \Pr\bigl(D_a f_\ell(x)=0\bigr)
  =\frac{r-1}{K-1}
\]
in the distinct-block case.  If $I_x$ denotes the indicator of this zero
derivative event, then
$\mathbb E[\sum_x I_x]=\sum_x\mathbb E[I_x]$; thus linearity of
expectation uses only the marginal probabilities and requires no
independence among the events indexed by $x$.  It now yields
\begin{equation}\label{eq:average}
  \frac{Q}{p}
  =T_a+(Q-T_a)\frac{r-1}{K-1}.
\end{equation}

Substituting $r=K/p$ into \eqref{eq:average} and multiplying by
$p(K-1)$ gives
\[
  Q(K-1)
  =pT_a(K-1)+(Q-T_a)(K-p).
\]
The right-hand side is
\[
  (p-1)KT_a+QK-pQ.
\]
After cancellation, we obtain
\[
  (p-1)Q=(p-1)KT_a.
\]
Since $p\ge2$, this proves the exact identity
\begin{equation}\label{eq:key-divisibility}
  KT_a=Q=p^n.
\end{equation}

Because $T_a$ is an integer, \eqref{eq:key-divisibility} implies
$K\mid p^n$.  Every positive divisor of the prime power $p^n$ is
$p^k$ for a unique integer $k$ with $0\le k\le n$.  Finally, the
definition of a bent partition includes $p\mid K$, so $k\ge1$.
Therefore $K=p^k$ with $1\le k\le n$, as claimed.
\end{proof}

\begin{remark}
The argument in fact uses only $n>0$, not the parity of $n$.  Thus the
same conclusion holds whenever the stated Walsh-flat bent functions and
the stated bent-partition property exist in odd dimension as well.
Moreover, the proof gives the exact identity
\[
  \#\bigl\{x\in V:\text{$x$ and $x+a$ lie in the same block}\bigr\}
  =\frac{p^n}{K}
\]
for every nonzero $a\in V$.
\end{remark}

\end{document}
